\documentclass[11pt]{article}
\usepackage[margin=1.15in]{geometry}
\usepackage{amsmath,amssymb,amsthm}
\usepackage{booktabs}
\usepackage{microtype}
\usepackage[hidelinks]{hyperref}
\usepackage{array}

\newtheorem{proposition}{Proposition}
\newtheorem{remark}{Remark}
\newcommand{\rr}{\rho}
\newcommand{\ts}{L_{T^*}}
\newcommand{\snet}{L_{S^+}\!-\!L_{S^-}}

\title{An Independent Verification of the Gilbert--Pollak Certificate
Layer and a Certificate-Layer Bound of $\rho \ge 0.860$ for the Steiner
Ratio}

\author{Claude Fable 5 (Anthropic)\thanks{Claude Fable 5 is an AI system.
This work was commissioned, directed, resourced, and approved for release
by J.~Savva (\texttt{aimathpapers@gmail.com}), who is the accountable
correspondent for all claims herein; see the AI Authorship Statement
(Section~\ref{sec:ai}). Code, data manifests, and reproduction
instructions: \url{https://github.com/aimathpapers/Steiner_Ratio},
archived at \url{https://doi.org/10.5281/zenodo.22223485}.}}

\date{August 2026}

\begin{document}
\maketitle

\begin{abstract}
The Gilbert--Pollak conjecture (1968) asserts that the Steiner ratio of
the Euclidean plane is $\sqrt{3}/2 \approx 0.8660254$. A 2026 preprint of
Ke et al.\ (arXiv:2601.22365) established the record lower bound
$\rr \ge 0.8559$ by a computer-assisted argument whose final layer is a
dataset of $330{,}193{,}755$ parameter boxes, each certifying
non-positivity of a designated splitting function. We report (i) the
first independent verification of that certificate layer --- a clean-room
re-implementation using ball arithmetic that confirms all
$330{,}193{,}755$ records with zero failures under semantics strictly
stronger than the published checker's, closing $28{,}845{,}102$ regions
over their entire boxes; (ii) an exact ceiling analysis: the splitting
functions are affine in $\rr$ and no side condition of the pipeline
depends on $\rr$, so each record's supremal certifiable $\rr$ is
computable in one pass, giving the published dataset the exact ceiling
$0.8559000606$ --- the bound was pinned by the generator's subdivision
tolerance, not by the mathematics; and (iii) a new record: regenerating
the partitions with the published generator retargeted to $0.860$ and
verifying every one of the resulting $1{,}104{,}177{,}103$ regions at the
exact rational $43/50$, with zero defects, yields the certificate-layer
bound $\rr \ge 0.860$, subject to the same (rho-uniform) lemma-layer
obligations as the published proof. Sampled analyses indicate the
published lemma set extends to at least $0.858$ on all case families and
to $0.866$ --- within $2.5\times10^{-5}$ of the conjectured optimum ---
on the d\_regular family. All software, manifests, and run records are
public; upstream artifacts are identified by pinned SHA-256 hashes.
\end{abstract}

\section{Introduction}

For a finite set $P$ of points in the Euclidean plane, a \emph{Steiner
minimal tree} (SMT) is a shortest network connecting $P$, where the
network may introduce auxiliary vertices; a \emph{minimum spanning tree}
(MST) is a shortest network whose vertices are exactly $P$. The
\emph{Steiner ratio} is
\[
\rr \;=\; \inf_{P}\; \frac{\mathrm{SMT}(P)}{\mathrm{MST}(P)} .
\]
Gilbert and Pollak \cite{GP68} conjectured $\rr = \sqrt{3}/2$, the value
attained by the vertices of an equilateral triangle. Lower bounds
progressed from $1/\sqrt{3}$ \cite{GH76} through $0.743$ \cite{CH78} and
$0.8$ \cite{DH83} to the Chung--Graham bound $0.824$ \cite{CG85}. Du and
Hwang announced a complete proof in 1990 \cite{DH90,DH92}; the argument's
central reduction was later found to be unsupported, and the conjecture
is regarded as open \cite{IT12}. The Chung--Graham bound stood for four
decades until Ke et al.\ \cite{PKU26} established $\rr \ge 0.8559$ in
2026 by a case analysis whose leaf obligations are discharged by
computer: a generator subdivides the parameter space of each local
configuration into boxes and assigns each box an inequality certificate,
and a checker verifies each certificate numerically.

Computer-assisted proofs of this shape have two distinguishable layers.
The \emph{certificate layer} is the finite dataset of per-box claims
together with the checking procedure --- objectively re-verifiable by
anyone. The \emph{lemma layer} is the human mathematics that justifies
the checking rules: the geometric lemmas behind each certificate template,
the correctness of the case decomposition, and (for checkers that sample
only box corners) monotonicity or unimodality claims about the certified
functions. For \cite{PKU26} the lemma layer rests in part on
computer-algebra sessions that were not published. This paper concerns
the certificate layer; every statement below of the form ``$\rr \ge x$''
is to be read as conditional on the lemma layer of \cite{PKU26}, whose
obligations we show transfer unchanged from $0.8559$ to any target
(Remark~\ref{rem:uniform}).

\paragraph{Contributions.}
\begin{enumerate}
  \item \textbf{Independent verification} (\S\ref{sec:verifier},
  \S\ref{sec:audit}). A clean-room verifier --- different language,
  different arithmetic (Arb ball arithmetic \cite{Arb17} via
  python-flint), no shared code --- confirms all $330{,}193{,}755$
  published records with zero certificate-layer failures, under
  strictly stronger semantics: all $2^n$ box vertices are checked
  (the published checker, we find, visits $2^{n-1}$), and
  $28{,}845{,}102$ regions are certified over their entire box,
  removing those regions' dependence on the unimodality claim. To our
  knowledge this is the first independent replication of the 0.8559
  certificate. Findings of the audit, none affecting validity, are
  reported in \S\ref{sec:findings}.
  \item \textbf{Exact ceilings by a one-pass critical-$\rr$ analysis}
  (\S\ref{sec:ceiling}). Each certified splitting function has the form
  $F(\rr) = \rr\,\ts + (\snet)$ with $\ts \ge 0$, and no other component
  of the pipeline reads $\rr$. Hence each record has a critical value
  $\rr^\ast = \min_{\text{vertices}} (L_{S^-}\!-\!L_{S^+})/\ts$,
  computable with certified enclosures, and the dataset's supremal
  certifiable target is the minimum of $\rr^\ast$ over all records ---
  no bisection, no re-verification per probe. The published dataset's
  exact ceiling is $\mathbf{0.8559000606}$: the generator subdivides
  each box only until its certificate clears zero by $10^{-6}$, so the
  published partition certifies its target and essentially nothing more.
  \item \textbf{The record} (\S\ref{sec:record}). Re-running the
  published generator with its target constant set to $0.860$ (a
  one-line change against the pinned upstream commit) and verifying all
  $1{,}104{,}177{,}103$ regenerated regions at the exact rational
  $43/50$ yields $\rr \ge 0.860$ at the certificate layer, with zero
  failures, zero invalid records, and zero inconclusive enclosures;
  $244{,}138{,}895$ regions are certified over their entire boxes.
  Sampled scouting places the lemma set's reach at $\ge 0.858$ for
  every case family and at $0.866$ for d\_regular.
\end{enumerate}

All code, tests, content-addressed manifests, and per-run records are
available at \url{https://github.com/aimathpapers/Steiner_Ratio} and
archived at \url{https://doi.org/10.5281/zenodo.22223485}. The
upstream repository and dataset carry no license and are not
redistributed; reproduction fetches them from their public locations and
verifies them against pinned SHA-256 hashes.

\section{The certificate framework of Ke et al.}\label{sec:framework}

We summarize the structure of \cite{PKU26} as far as needed; the
notation follows their artifact. The proof reduces the conjecture-bound
question to four \emph{case families} --- two geometric configurations
(\texttt{d\_regular}, \texttt{d\_steiner}), each split by the boundary
condition $f \ge d$ or $f \le d$ on two of its parameters. A local
configuration is described by $n$ parameters ($n = 5,6,7,8$ per case),
so each case's domain is a box in $\mathbb{R}^n$, with some coordinates
unbounded above.

For each case the authors fix a finite menu of \emph{splittings}
(combinatorial ways a hypothetical worst-case network could be cut and
reconnected) and a menu of \emph{lemmas} (human-proved geometric bounds
on lengths of trapped components). For a splitting $s$ and lemma $\ell$,
the \emph{splitting function} has the form
\begin{equation}\label{eq:F}
F_{s,\ell}(\rr; x) \;=\; \rr \cdot \ts(x) \;+\; L_{S^+}(x) \;-\;
L_{S^-}(x),
\end{equation}
where $\ts$ is a sum of edge lengths and lemma bounds (each
non-negative), $L_{S^+}$ prices the reconnection (a distance, a
three-point Steiner length, or a lemma bound), and $L_{S^-}$ sums deleted
edge lengths. Non-positivity of $F_{s,\ell}$ over a parameter box
certifies that no configuration in that box can violate the bound $\rr$.
Lemma applicability is gated by side conditions on $x$ (never on $\rr$).

A generator program subdivides each case's domain into boxes and, for
each box, records a splitting and lemma under which non-positivity holds;
the published dataset comprises $330{,}193{,}755$ such records
(Table~\ref{tab:audit}). The published checker verifies each record by
evaluating $F$ at box corner points in floating-point interval
arithmetic; corner sampling is justified by an unpublished axis-wise
unimodality claim (their ``Shape Constraint''). Records whose box
violates the subcase boundary condition are skipped by rule, as in the
published checker.

\section{An independent verifier}\label{sec:verifier}

Our verifier shares no code, language, or numerical library with the
published stack (Julia / IntervalArithmetic.jl). It is written in Python
over Arb \cite{Arb17} ball arithmetic (python-flint): every quantity is
carried as a rigorous enclosure, and a comparison is accepted only when
it holds for \emph{every} point of the enclosures. Three design points
matter for what follows.

\textbf{Exact $\rr$.} The target is held as an exact rational
($8559/10000$, later $43/50$) and converted to a ball only at evaluation,
at the working precision. No floating-point representation of the claim
enters any decision.

\textbf{Stronger vertex semantics.} We evaluate $F$ at all $2^n$ finite
vertices of each box (with unbounded directions handled exactly as the
published rules prescribe), escalating the working precision through a
$128/256/512$-bit ladder before recording anything but a clean pass, so
enclosure width is never mistaken for a failure.

\textbf{Full-box certificates.} Where possible we additionally certify
$F \le 0$ over the \emph{entire} box by adaptive bisection with a bounded
budget. A region so certified no longer depends on the unimodality claim
at all.

The verifier is tested at two seams: a certificate-boundary suite over
hand-constructed records with known outcomes, and an oracle suite that
compares every lemma bound, condition, and composed $F$ against an
independently written 60-digit mpmath reference at exact dyadic points
(containment of the reference value in our enclosure, and agreement of
branch decisions away from boundaries). The reference implementation is
also embedded in self-contained ``capsule'' replay scripts that would
accompany any defect claim; none were needed.

\section{Audit of the published certificate}\label{sec:audit}

\begin{table}[t]
\centering
\begin{tabular}{lrrrr}
\toprule
case & records & pass (vertex) & pass (full box) & skipped \\
\midrule
d\_regular, $f\ge d$ & 1{,}888{,}501 & \multicolumn{3}{c}{\emph{all records verified; zero failures}}\\
d\_regular, $f\le d$ & 63{,}466{,}354 & & & \\
d\_steiner, $f\ge d$ & 5{,}402{,}950 & & & \\
d\_steiner, $f\le d$ & 259{,}435{,}950 & & & \\
\midrule
total & 330{,}193{,}755 & \multicolumn{2}{r}{265{,}761{,}575 checked} & 64{,}432{,}180 \\
\bottomrule
\end{tabular}
\caption{The published dataset (target $0.8559$), as verified by our
stack. ``Skipped'' records are excluded by the published subcase rules
(their checker would not evaluate them either); every skip was
re-derived, not trusted. $28{,}845{,}102$ of the checked regions were
additionally certified over their entire box.}
\label{tab:audit}
\end{table}

Every record of the published dataset verifies: no failures, no invalid
records, no inconclusive enclosures at any precision rung
(Table~\ref{tab:audit}). We regard the certificate layer of \cite{PKU26}
as independently confirmed --- indeed strengthened, since our semantics
subsume the published checker's.

\subsection{Findings}\label{sec:findings}

The audit surfaced five findings, none affecting the validity of the
published certificate, all relevant to anyone building on it.

\begin{enumerate}
  \item[F1.] Regenerating the dataset from the published generator
  source reproduces the mathematical content exactly (region bounds and
  identities), while canonical orderings and platform-dependent float
  formatting differ --- regeneration is deterministic in content, not in
  bytes.
  \item[F2.] The published checker's vertex loop indexes mask bits
  $1..n$, visiting $2^{n-1}$ of the $2^n$ vertex assignments it
  documents. Our all-vertices verification shows the stronger property
  holds anyway.
  \item[F3.] Sentinel encodings differ between the published data
  (records excluded by subcase rules carry splitting id $0$) and the
  current generator source (which emits $-1$): the published dataset was
  not produced by the repository's HEAD.
  \item[F4.] No upstream artifact was content-addressed; we published
  pinned SHA-256 manifests for the repository snapshot (commit
  \texttt{709673a8}) and every dataset file.
  \item[F5.] The v2 paper describes a Mathematica/CAD verification and
  lists interval arithmetic as future work, while the shipped artifact
  is a Julia interval checker --- the peer-reviewable text and the actual
  proof artifact have diverged.
\end{enumerate}

\section{Critical $\rr$ and exact ceilings}\label{sec:ceiling}

\begin{proposition}\label{prop:affine}
Fix a record (a box $B$, splitting $s$, lemma $\ell$) admitted by the
side conditions. As a function of the target, $F_{s,\ell}(\rr; x) = \rr
\cdot \ts(x) + (\snet)(x)$ with $\ts(x) \ge 0$, and neither the side
conditions nor the skip rules depend on $\rr$. Consequently, for each
checked vertex $x$ the set of targets certified at $x$ is an interval
$(-\infty, \rr^\ast(x)]$ (or all targets, when $\ts(x)=0$ and
$(\snet)(x) \le 0$), with
\[
\rr^\ast(x) \;=\; \frac{L_{S^-}(x) - L_{S^+}(x)}{\ts(x)},
\]
and the record's supremal certifiable target is $\rr^\ast(B) =
\min_{x \in V(B)} \rr^\ast(x)$ over the checked vertex set $V(B)$. The
dataset's supremal certifiable target is the minimum of $\rr^\ast(B)$
over its records.
\end{proposition}

\begin{proof}
Affineness and the sign of $\ts$ are immediate from
\eqref{eq:F}: $\ts$ is a sum of Euclidean lengths and lemma
bounds, each non-negative on its admissible domain. Monotonicity of
$\rr \mapsto F$ in $\rr$ then gives the interval structure at each
vertex; intersecting over vertices and over records gives the minima.
Since side conditions and skip rules are $\rr$-free, the checked vertex
set $V(B)$ and the admitted record set are the same for every target,
so the minima are well-defined independently of $\rr$.
\end{proof}

Ball arithmetic delivers certified two-sided enclosures of each
$\rr^\ast(B)$ (division of enclosures, with zero-straddling denominators
handled conservatively: a certified $(\snet)(x) > 0$ means the vertex
fails at every target since $\ts \ge 0$; a denominator enclosure touching
zero with $(\snet)(x) \le 0$ still yields a valid lower bound). One pass
over the dataset therefore computes the exact ceiling. For full-box
certificates the same ratio, evaluated on each leaf of the certifying
subdivision, gives a one-sided bound valid over the whole box.

\begin{table}[t]
\centering
\begin{tabular}{lrl}
\toprule
case & records & exact ceiling of the published partition \\
\midrule
d\_regular, $f\ge d$ & 1{,}888{,}501 & $0.8559001279$ \\
d\_regular, $f\le d$ & 63{,}466{,}354 & $0.8559000619$ \\
d\_steiner, $f\ge d$ & 5{,}402{,}950 & $0.8559000943$ \\
d\_steiner, $f\le d$ & 259{,}435{,}950 & $\mathbf{0.8559000606}$ \\
\bottomrule
\end{tabular}
\caption{Exact supremal certifiable targets of the published dataset
(certified enclosure widths below $10^{-10}$; zero inconclusive
records). The published certificate proves $\rr \ge 0.8559$ and nothing
materially more.}
\label{tab:ceilings}
\end{table}

Table~\ref{tab:ceilings} gives the result: the published dataset's
ceiling is $0.8559000606$. The explanation is a single line of the
generator: each box is subdivided only until its certificate clears zero
by a tolerance of $10^{-6}$, so the partition is tailored to its target
with essentially no slack. The bound $0.8559$ was pinned by a stopping
rule. This immediately suggests the experiment of the next section.

\begin{remark}\label{rem:uniform}
Every lemma-layer obligation inherited from \cite{PKU26} --- the
geometric lemmas, the case decomposition, monotonicity of constituent
lengths, and the unimodality claim where vertex checking is used --- is a
statement about lengths and geometry in which $\rr$ does not appear.
The obligations therefore transfer unchanged to any target; certifying a
regenerated partition at a higher $\rr$ incurs no new lemma-layer debt.
\end{remark}

\section{Regeneration and the record}\label{sec:record}

\begin{table}[t]
\centering
\begin{tabular}{lrrrrr}
\toprule
case & records & vertex pass & full box & skipped & defects \\
\midrule
d\_regular, $f\ge d$ & 105{,}508{,}647 & 88{,}464{,}876 & 17{,}043{,}771 & 0 & 0 \\
d\_regular, $f\le d$ & 306{,}908{,}065 & 189{,}196{,}066 & 97{,}215{,}618 & 20{,}496{,}381 & 0 \\
d\_steiner, $f\le d$ & 291{,}063{,}158 & 231{,}250{,}637 & 46{,}280{,}442 & 13{,}532{,}079 & 0 \\
d\_steiner, $f\ge d$ & 400{,}697{,}233 & 316{,}783{,}555 & 83{,}599{,}064 & 314{,}614 & 0 \\
\midrule
total & 1{,}104{,}177{,}103 & & 244{,}138{,}895 & & \textbf{0} \\
\bottomrule
\end{tabular}
\caption{The record verification: partitions regenerated at target
$0.860$, every record verified at the exact rational $43/50$ under
full-box-first semantics. SHA-256 hashes of the four certificate files
are recorded in the artifact manifests.}
\label{tab:record}
\end{table}

We re-ran the published generator --- from the pinned snapshot, in
isolated working directories, with exactly one source change per case
family: the constant \texttt{rho = 0.8559} set to the new target --- and
verified the regenerated partitions with the stack of
\S\ref{sec:verifier} at the exact rational target. As a first step,
$\rr = 0.856$ was certified exhaustively ($559{,}929{,}762$ regions,
zero defects). The main result is the target $0.860$:
all $1{,}104{,}177{,}103$ regenerated regions verify at $43/50$ with
zero failures, zero invalid records, and zero inconclusive enclosures
(Table~\ref{tab:record}); $244{,}138{,}895$ regions carry full-box
certificates. With Remark~\ref{rem:uniform}:

\begin{quote}
\textbf{Certificate-layer bound.} Subject to the lemma layer of
\cite{PKU26}, the Steiner ratio satisfies $\rr \ge 0.860$.
\end{quote}

\paragraph{Cost.} Regeneration took between $0.9$ and $9.6$ hours per
case (single-threaded, Apple M-series); partition sizes grew by factors
between $1.1$ and $32$ relative to the published ones. Verification
throughput was $\approx 4{,}600$ records/second at 16 worker processes;
the largest case verified in about two days.

\paragraph{Scouting above 0.860.} Sampled critical-$\rr$ analyses
(50{,}000 records per partition, certified per-record enclosures, zero
adverse outcomes in every sample) were run on regenerated partitions at
higher targets: for d\_steiner, targets $0.857$ and $0.858$; for
d\_regular, targets through $0.866$, where the sampled partition still
certifies its target --- within $2.5 \times 10^{-5}$ of $\sqrt{3}/2$.
These are scouting results on samples, not claims; we state as claims
only the exhaustively verified targets $0.856$ and $0.860$.

\section{Scope and threats to validity}\label{sec:scope}

\textbf{The lemma layer is inherited, not verified.} Our result is
exactly as conditional as the published one, at a higher number. The
specific obligations: the geometric lemmas behind certificate templates
(partly unpublished computer-algebra sessions), the completeness of the
case decomposition, and, for vertex-checked regions only, the axis
unimodality claim. Full-box regions ($244$M of $1.1$B at $0.860$) do
not use unimodality.

\textbf{Partition coverage is the generator's construction.} We verify
that every emitted region certifies its inequality; that the emitted
regions cover each case domain is a property of the generator's
recursive subdivision, exactly as in the published work. Our audit
supports (content-identical regeneration, F1) but does not prove
generator correctness.

\textbf{Trust in the toolchain.} Our verifier trusts Arb's ball
arithmetic and the Python runtime; the oracle seam (an independent
60-digit implementation) and the two-seam test suite (128 tests) are the
mitigation. The verifier never trusts the generator's arithmetic:
certificates are re-derived from the record coordinates alone.

\textbf{What would falsify the claim.} A region of any regenerated
partition whose splitting function is positive somewhere in its box; a
demonstrated coverage gap in the generator's subdivision; or a
counterexample to an inherited lemma. The artifact includes a
capsule mechanism producing a minimal, dependency-light replay script
for any alleged defect.

\section{Future directions}\label{sec:future}

The scouting data indicate the published lemma set is close to exhausted
somewhere above $0.86$; further progress will require new lemmas ---
new local geometry, rigorously certified, for which our infrastructure
(bottleneck extraction, candidate falsification, replayable lemma
certificates) is built. No finite refinement reaches $\sqrt{3}/2$: the
conjecture's endgame requires a global-structure argument of exactly the
kind that failed in 1990 \cite{DH92,IT12}, and the certificate program
should be seen as rigorously draining the neighborhood of that hard
core. A machine-checked formalization (e.g.\ in Lean) of both layers is
the natural terminal state for a problem with this history; the
artifact's exact-rational, replayable structure is designed to make that
translation tractable.

\section{AI Authorship Statement}\label{sec:ai}

This paper, and the work it reports --- the verifier implementation, the
test suites, the audit, the critical-$\rr$ analysis, the regeneration
campaign, and the drafting of this manuscript --- were carried out by
Claude Fable 5, an AI system built by Anthropic, operating the
computation autonomously over approximately ten days. J.~Savva
commissioned and directed the work, set its scope, rigor bar, and stopping
decisions, provided the computing resources, reviewed the outputs, and is
solely accountable for the decision to publish and for the claims made
here. No mathematical statement in this paper rests on the AI's
judgment: every verification outcome is the result of certified interval
arithmetic, every artifact is content-addressed, and every claim is
reproducible from the public repository by the commands documented
there. The mathematical machinery being verified and extended is due to
the authors of \cite{PKU26}.

\subsection*{Acknowledgments}
Ke, Huang, Shu, He, Gai, and Wang \cite{PKU26} built the lemma set, case
decomposition, and generator on which everything here stands; this work
would not exist without theirs. F.~Johansson's Arb library \cite{Arb17}
provides the certified arithmetic.

\begin{thebibliography}{9}
\bibitem{GP68} E.~N.~Gilbert and H.~O.~Pollak, Steiner minimal trees,
\emph{SIAM J. Appl. Math.} 16 (1968), 1--29.
\bibitem{GH76} R.~L.~Graham and F.~K.~Hwang, Remarks on Steiner minimal
trees, I, \emph{Bull. Inst. Math. Acad. Sinica} 4 (1976), 177--182.
\bibitem{CH78} F.~R.~K.~Chung and F.~K.~Hwang, A lower bound for the
Steiner tree problem, \emph{SIAM J. Appl. Math.} 34 (1978), 27--36.
\bibitem{DH83} D.-Z.~Du and F.~K.~Hwang, A new bound for the Steiner
ratio, \emph{Trans. Amer. Math. Soc.} 278 (1983), 137--148.
\bibitem{CG85} F.~R.~K.~Chung and R.~L.~Graham, A new bound for
Euclidean Steiner minimal trees, \emph{Ann. N.Y. Acad. Sci.} 440 (1985),
328--346.
\bibitem{DH90} D.-Z.~Du and F.~K.~Hwang, The Steiner ratio conjecture
of Gilbert and Pollak is true, \emph{Proc. Natl. Acad. Sci. USA} 87
(1990), 9464--9466.
\bibitem{DH92} D.-Z.~Du and F.~K.~Hwang, A proof of the Gilbert--Pollak
conjecture on the Steiner ratio, \emph{Algorithmica} 7 (1992), 121--135.
\bibitem{IT12} A.~O.~Ivanov and A.~A.~Tuzhilin, The Steiner ratio
Gilbert--Pollak conjecture is still open, \emph{Algorithmica} 62 (2012),
630--632.
\bibitem{PKU26} Y.~Ke, T.~Huang, Y.~Shu, D.~He, J.~Gai, and L.~Wang,
Towards solving the Gilbert--Pollak conjecture via large language
models, arXiv:2601.22365 (2026), v2 May 2026,
DOI 10.48550/arXiv.2601.22365; artifact:
\texttt{github.com/keyisi2006/Steiner-Ratio}, pinned at commit
\texttt{709673a8}.
\bibitem{Arb17} F.~Johansson, Arb: efficient arbitrary-precision
midpoint-radius interval arithmetic, \emph{IEEE Trans. Computers} 66(8)
(2017), 1281--1292.
\end{thebibliography}

\end{document}
